% Il metodo: le fasi in avanti e i ritorni che hanno cambiato il progetto.
\begin{tikzpicture}[
  font=\footnotesize, >=Stealth,
  fase/.style={draw=fblu, fill=fblu!7, rounded corners=2pt, align=center, text width=18mm,
               minimum height=14mm, inner sep=2pt},
  ritorno/.style={->, farancio, thick, rounded corners=3pt},
  lab/.style={font=\scriptsize, farancio, fill=white, inner sep=1.5pt}]
  \node[fase] (s) at (0,0)     {specification\\and criteria\\[1pt]\scriptsize Ch.~\ref{ch:spec}};
  \node[fase] (a) at (2.2,0)   {analytical\\sizing\\[1pt]\scriptsize Ch.~\ref{ch:analytical}--\ref{ch:modulation}};
  \node[fase] (m) at (4.4,0)   {circuit models\\{\scriptsize PLECS, SPICE}\\[1pt]\scriptsize Ch.~\ref{ch:model}--\ref{ch:losses}};
  \node[fase] (t) at (6.6,0)   {thermal\\and cooling\\[1pt]\scriptsize Ch.~\ref{ch:thermal}};
  \node[fase] (x) at (8.8,0)   {auxiliary\\circuits\\[1pt]\scriptsize Ch.~\ref{ch:aux}};
  \node[fase] (b) at (11.0,0)  {schematic\\and board\\[1pt]\scriptsize Ch.~\ref{ch:schematic}--\ref{ch:layout}};
  \node[fase, draw=fgrigio, fill=fgrigio!10] (v) at (13.2,0) {bench plan\\[1pt]\scriptsize App.~\ref{app:assumptions}};
  \foreach \i/\j in {s/a, a/m, m/t, t/x, x/b, b/v} \draw[->, thick, fblu] (\i) -- (\j);
  % ritorni: ciascuno e' una scoperta che ha cambiato una decisione a monte
  \draw[ritorno] (6.6,-0.7) -- (6.6,-1.15) -- node[lab]{locked rotor changes the device} (0,-1.15) -- (0,-0.7);
  \draw[ritorno] (11.0,-0.7) -- (11.0,-1.75) -- node[lab]{drawn rail replaces the ideal plate} (6.9,-1.75) -- (6.9,-0.7);
  \draw[ritorno] (4.1,0.7) -- (4.1,1.15) -- node[lab]{C2 sets $R_g$ and dead time} (0,1.15) -- (0,0.7);
  \draw[ritorno] (11.0,0.7) -- (11.0,1.75) -- node[lab]{extracted loop re-simulated against C2} (4.7,1.75) -- (4.7,0.7);
\end{tikzpicture}
